\section{Архитектура модели-оценщика}
\label{sec:model}

\subsection{Общая схема}

Модель-оценщик представляет собой базовую трансформерную модель
с тремя классификационными головами на уровне токенов, соответствующими
целям TRACe: relevance, utilization, adherence. На вход модели поступает
объединённая последовательность $[q][\text{SEP}][d_1, \dots, d_{|D|}][\text{SEP}][a]$,
токенизированная единым токенизатором \glslink{backbone}{базовой модели}. На всех слоях
используется стандартный механизм self-attention; дополнительные межтокенные зависимости в
 классификационных головах не моделируются.

\begin{figure}[H]
\centering
\begin{tikzpicture}[
  node distance=7mm and 14mm,
  block/.style={draw, rounded corners=2pt, align=center, minimum height=8mm, inner sep=2pt, font=\small},
  head/.style={draw, rounded corners=2pt, align=center, fill=gray!10, minimum height=8mm, font=\small},
  arr/.style={-latex}
]
\node[block] (input) {$[q]\,[\text{SEP}]\,[d_1,\dots,d_{|D|}]\,[\text{SEP}]\,[a]$};
\node[block, below=of input] (tok) {Токенизатор};
\node[block, below=of tok, fill=blue!8] (bb) {Базовая модель (DeBERTa / BGE-M3 / RuModernBERT)};
\node[block, below=of bb] (hs) {Hidden states $H \in \mathbb{R}^{T \times d}$};
\node[head, below left=8mm and 5mm of hs] (rel)  {Relevance\\ head};
\node[head, below=8mm of hs] (util) {Utilization\\ head};
\node[head, below right=8mm and 5mm of hs] (adh)  {Adherence\\ head};
\node[block, below=22mm of hs, align=center] (loss) {$\mathcal{L} = \dfrac{\alpha_R\mathcal{L}_R + \alpha_U\mathcal{L}_U + \alpha_A\mathcal{L}_A}{\alpha_R + \alpha_U + \alpha_A}$\\ (BCE по валидным токенам)};

\draw[arr] (input) -- (tok);
\draw[arr] (tok) -- (bb);
\draw[arr] (bb) -- (hs);
\draw[arr] (hs) -- (rel);
\draw[arr] (hs) -- (util);
\draw[arr] (hs) -- (adh);
\draw[arr] (rel) -- (loss);
\draw[arr] (util) -- (loss);
\draw[arr] (adh) -- (loss);
\end{tikzpicture}
\caption{Общая схема токен-уровневого оценщика TRACe.}
\label{fig:model_overview}
\end{figure}

\subsection{Две архитектуры головы}
\label{sec:heads}

\paragraph{Простая голова (Simple head).} Воспроизводит архитектуру из оригинальной статьи
 RAGBench~\cite{friel2024ragbench}: после энкодера применяется общий слой
  \texttt{Dropout} \cite{srivastava2014dropout} и три независимых линейных слоя $\mathbb{R}^d \to \mathbb{R}^1$, 
  по одному для каждого таргета. 
Формально, для скрытого состояния $h_t$ токена $t$:
\[
\ell^{(k)}_t = W^{(k)} h_t + b^{(k)}, \quad k \in \{R, U, A\},
\]
где $W^{(k)} \in \mathbb{R}^{1 \times d}$. Число параметров в головах составляет $3(d+1)$ - порядка 2.3~тыс. для базовой модели с $d=768$ и 4.6~тыс. для $d=1024$.

\paragraph{Сложная голова (Complex head).} Архитектура с общим MLP-преобразованием \cite{MLP} и индивидуальными 
головами-MLP с функцией активации GELU \cite{hendrycks2016gaussian}:
\begin{align*}
x_t &= \mathrm{LN}\bigl(\mathrm{Dropout}(\mathrm{GELU}(W_1 h_t + b_1))\bigr) \in \mathbb{R}^{d_s}, \\
\tilde{h}_t &= \mathrm{LN}\bigl(W_2 x_t + h_t\bigr),\\ 
\ell^{(k)}_t &= V_2^{(k)} \bigl(\mathrm{Dropout}(\mathrm{GELU}(\mathrm{LN}(V_1^{(k)} \tilde{h}_t)))\bigr).
\end{align*}
Здесь $d_s$ совпадает с $d$ (\texttt{shared\_mult}~=~1), а внутренняя размерность
 головы равна $d/2$ (\texttt{head\_mult}~=~0.5). Число параметров в головах возрастает
  примерно на два порядка: $\sim$2.5~M для $d=768$.

Сравнение двух вариантов вынесено в раздел~\ref{sec:ablation_head}. Здесь важно лишь
отметить, что обе архитектуры рассматриваются как конкурирующие гипотезы, а их
практическая устойчивость проверяется уже на этапе экспериментов.

\subsection{Функция потерь}

Лосс задаётся как нормированная взвешенная сумма бинарных кросс-энтропий по валидным токенам:
\[
\mathcal{L}(\theta) =
\frac{\alpha_R \mathcal{L}_R + \alpha_U \mathcal{L}_U + \alpha_A \mathcal{L}_A}
{\alpha_R + \alpha_U + \alpha_A},
\]
где для $k \in \{R, U\}$ суммирование идёт по маске документов $m^{(C)}_t$, а для $k = A$ -
 по маске ответа $m^{(A)}_t$. Рассматриваются две схемы весов:
\begin{itemize}
  \item $(\alpha_R, \alpha_U, \alpha_A) = (1, 1, 1)$ - базовая;
  \item $(\alpha_R, \alpha_U, \alpha_A) = (3, 3, 1)$ - гипотеза: adherence~
  семантически более простая задача (булев флаг по предложению, а не расстановка токенных меток),
   и её пониженный вес позволит сильнее сконцентрировать обучение на релевантности и использовании.
\end{itemize}
Далее в разделе~\ref{sec:ablation_weights} эта гипотеза проверяется экспериментально;
как будет показано, различие между двумя схемами весов оказывается небольшим.

\subsection{Подбор порогов}
\label{sec:thresholds}

После обучения каждая голова выдаёт \glslink{логит}{логиты} $\ell^{(k)}_t$,
которые превращаются в бинарные предсказания через сигмоид и порог:
 $\hat{y}^{(k)}_t = \mathbb{1}[\sigma(\ell^{(k)}_t) \ge \tau_k]$.

Пороги $\tau_R, \tau_U, \tau_A$ подбираются независимо для каждой головы 
на валидационной выборке. Реализация в \texttt{model\_utils/metrics.py} 
сводится к одному проходу по валидации для сбора вероятностей и перебору сетки
 кандидатов $\tau \in \{0{.}1, 0{.}2, \dots, 0{.}7\}$. Для каждого 
 кандидата вычисляется F1-мера на уровне токенов по маске соответствующего
  таргета, после чего выбирается порог, максимизирующий данную метрику.

Как показывают результаты раздела~\ref{sec:thresholds_results}, оптимальные
пороги заметно отклоняются от 0.5 и зависят как от модели, так и от таргета.
Для relevance и utilization медианный порог на русскоязычных конфигурациях
лежит около 0.4, с разбросом по диапазону $[0.2, 0.5]$. Для adherence оптимум
часто приходится на нижнюю границу сетки ($\tau_A \approx 0{.}1$), что
объясняется высокой долей положительных примеров: модель уже на низком пороге
выдаёт почти все положительные предсказания.

\subsection{Выбор базовой модели}
\label{sec:backbones}

В основной сетке экспериментов используются три базовые модели, охватывающие основные классы архитектур, применимых к русскоязычной задаче:

\begin{itemize}
  \item \textbf{DeBERTa-v3-large} (\texttt{MoritzLaurer/DeBERTa-v3-
  large-mnli-fever-anli-ling-wanli}) - англоязычная базовая конфигурация из оригинальной
  статьи RAGBench, предобучена на NLI-задачах. Максимальная длина контекста: 512.
   Используется как контрольный эксперимент и для сравнения RU/EN.
  \item \textbf{BGE-M3} (\texttt{BAAI/bge-m3} \cite{bge-m3} ) - мультиязычная модель,
  обученная для поиска; архитектурно - XLM-RoBERTa-large с поддержкой длинного контекста
   (до 8192 через масштабирование позиционных кодировок).
  \item \textbf{RuModernBERT} (\texttt{deepvk/RuModernBERT-base} \cite{deepvk2025rumodernbert})
     - русскоязычная модель ModernBERT, обученная на русскоязычном корпусе;
     поддерживает обработку последовательностей длиной до 8192 токенов за счёт использования
     вращательных позиционных эмбеддингов (rotary position embeddings \cite{su2024roformer}).
\end{itemize}

Дополнительно на предварительном этапе рассматривалась модель \textbf{ruRoPEBert-e5-base-2k}
 (\texttt{Tochka-AI/ruRoPEBert-e5-base-2k}) - русскоязычный BERT с RoPE и контекстом до 2048
  токенов, предобученный на задачах поиска. Результаты предварительного эксперимента
   (раздел~\ref{sec:rope_preliminary}) показали, что эта модель существенно уступает
    альтернативам в задаче токен-уровневой классификации, поэтому она была исключена из
    основной сетки.

Мотивация выбора: (а)~DeBERTa дублирует базовую конфигурацию статьи и даёт
верхнюю оценку качества на английском; (б)~BGE-M3 - один из популярных сильных
публично доступных мультиязычных энкодеров с длинным контекстом;
(в)~RuModernBERT - современная русскоязычная модель с rotary-позиционными
 кодировками и поддержкой длинного контекста до 8192 токенов. Таким образом, сетка охватывает три
 ортогональных оси: язык предобучения (EN / мультиязычная / RU), размер модели и схему
  позиционных кодировок (абсолютные vs rotary).

\subsection{Предварительный эксперимент с ruRoPEBert}
\label{sec:rope_preliminary}

Модель ruRoPEBert-e5-base-2k запускалась в идентичных с остальными базовыми моделями
 условиях (RU, Simple-голова, веса лосса 1:1:1) на длинах контекста 512, 1024 и 2048 как
  предварительная проверка перед включением в основную сетку. На всех трёх длинах F1 по
   relevance и utilization оставался практически нулевым; adherence при этом работал на
    уровне $\approx 0.85$, однако само по себе такое значение не свидетельствует о
     хорошем качестве модели, поскольку доля положительных примеров по adherence высока, и
      константный предсказатель <<всегда true>> достигает сопоставимого F1
      (см. обсуждение в разделе~\ref{sec:ablation_head}). По этой причине ruRoPEBert был
      исключён из основной сетки, чтобы сосредоточить вычислительный бюджет на более
      перспективных конфигурациях.

Возможное объяснение этого результата состоит в том, что предобучение ruRoPEBert по
 схеме e5 ориентировано прежде всего на задачу поиска (контрастивное обучение
 эмбеддингов предложений), что плохо согласуется с классификацией фрагментов
 релевантных и использованных предложений на уровне токенов. В отличие от DeBERTa
 (NLI-предобучение, более близкое к задаче следования) и RuModernBERT
 (универсальная MLM-модель), ruRoPEBert, по-видимому, формирует представления,
 лучше приспособленные к агрегированному сопоставлению, чем к позиционному
 разметочному сигналу. Помимо этого, нельзя полностью исключать и инженерные
 причины (например, особенности агрегации скрытых состояний или взаимодействия
 модели с маской на длинном входе); проверка таких гипотез потребовала бы
 отдельных запусков, и в условиях ограниченного вычислительного бюджета было
 принято решение перераспределить ресурсы в пользу более перспективных
 базовых моделей. Поэтому приведённая интерпретация остаётся рабочей гипотезой,
 хотя и хорошо согласуется с наблюдаемыми результатами.

\begin{table}[H]
\centering
\small
\caption{Результаты предварительного эксперимента на длине контекста 512 (RU, Simple,
 веса 1:1:1). ruRoPEBert выдаёт нулевой F1 на relevance и utilization; значение
  adherence $\approx 0.85$ объясняется более легким предсказанием, так как этот таргет всегда в конце.}
\label{tab:rope_preliminary_512}
\begin{tabular}{lccc}
\toprule
Модель & $F_1^{(R)}$ & $F_1^{(U)}$ & $F_1^{(A)}$ \\
\midrule
DeBERTa-v3-large        & 0.73 & 0.75 & 0.897 \\
RuModernBERT-base       & 0.71 & 0.74 & 0.869 \\
ruRoPEBert-e5-base-2k   & 0.00 & 0.00 & 0.852 \\
\bottomrule
\end{tabular}
\end{table}

Зависимость F1 от длины контекста с разбиением по таргетам для всех трёх моделей
 приведена на рис.~\ref{fig:rope_preliminary}: у ruRoPEBert кривые для relevance
  и utilization остаются на нуле во всём проверенном диапазоне длин, тогда как у
   DeBERTa и RuModernBERT значения лежат в рабочем диапазоне $0.7\text{--}0.9$.

\begin{figure}[H]
\centering
\includegraphics[width=0.95\textwidth]{rope_preliminary_f1_vs_length.png}
\caption{Зависимость $F_1$ от длины контекста с разбиением по таргетам (relevance,
 utilization, adherence) для DeBERTa-v3-large, RuModernBERT-base и ruRoPEBert-e5-base-2k
  (RU, Simple, веса 1:1:1). У ruRoPEBert линии для relevance и utilization лежат на нуле
   при всех длинах.}
\label{fig:rope_preliminary}
\end{figure}

\subsection{Детали обучения}
\label{sec:training_details}

Все эксперименты выполнялись в единой реализации 
(модуль \texttt{model\_utils/}, конвейер \texttt{runner.py} + \texttt{train\_loop.py}):
\begin{itemize}
  \item \textbf{Оптимизатор}: AdamW с раздельными группами параметров для 
  базовой модели ($\mathrm{lr} = 2{\times}10^{-5}$) и голов 
  ($\mathrm{lr} = 3{\times}10^{-4}$), \texttt{weight\_decay}~=~0.001.
  \item \textbf{Warmup}: линейный разогрев на 10\% шагов, затем линейное затухание до нуля.
  \item \textbf{Батч}: 64 для обучения, 32 для валидации/теста. Для длин~$\ge$2048 батч
   уменьшался до 2-4 с gradient accumulation.
  \item \textbf{Точность}: FP16 mixed precision (torch AMP); \texttt{grad\_clip}~=~1.0.
  \item \textbf{Число эпох}: 4 для всех конфигураций; при этом сохранялась лучшая по 
  качеству на валидационной выборке модель, что позволяло ограничить переобучение.
\end{itemize}

Подбор гиперпараметров выполнялся только на 
оси <<архитектура $\times$ длина $\times$ веса лосса>>: 
остальные параметры (lr, warmup, dropout, weight decay) зафиксированы.
 Такой дизайн позволяет интерпретировать вариации F1 как эффект изучаемых факторов, 
 а не случайный шум гиперпараметрического поиска.
